% Volume IX, printed pp. 112--121.
\begingroup
\clearpage
\noindent112\hfill\textsc{Posthumous Papers.}\hfill\mbox{}\par\bigskip
\begin{center}[6.]\end{center}
For the computation [of $\psi$ or $\varphi$], one may also use the following when
\[
\varphi=P+\delta
\]
is set:
\[
\begin{aligned}
\psi&=P+\frac{1-ee}{1-ee\sin P^2}\delta+\tfrac34\frac{ee(1-ee)\sin2P}{(1-ee\sin P^2)^2}\delta\delta+\text{ etc.}\\[5pt]
&=\varphi-\frac{ee\cos P^2}{1-ee\sin P^2}\delta+\tfrac34\frac{ee(1-ee)\sin2P}{(1-ee\sin P^2)^2}\delta\delta+\text{ etc.};
\end{aligned}
\]
[when]
\[
\psi=P+\varepsilon
\]
is set:
\[
\varphi=\psi+\frac{ee\cos P^2}{1-ee\sin P^2}\varepsilon-\tfrac34\frac{ee(1-ee\sin P^2)\sin2P}{(1-ee)^2}\varepsilon\varepsilon-\text{ etc.}
\]
\bigskip
\begin{center}[7.]\end{center}
[If one sets]
\[
\begin{aligned}
e\sin\varphi&=\sin u\\
e\sin P&=\sin U,
\end{aligned}
\]
[then one also has]
\[
\psi=\varphi-\cos\tfrac12(\varphi+\psi).Ae\log\frac{\operatorname{tang}(45^\circ+\frac12u)}{\operatorname{tang}(45^\circ+\frac12U)}\ldots,
\]
where
\[
A=\frac{206\,265}{\mathrm{Mod.}}.
\]
[For a flattening of $\dfrac1{302{,}68}$,] $\log Ae=4{,}586\,3052$.

To a close approximation [we have]
\[
\varphi=\psi+\frac{ee}{1-ee}(\psi-P)\cos(\tfrac14P+\tfrac34\psi)^2,
\]
more accurately
\[
\varphi=\psi+\frac{ee}{1-ee}(\psi-P)\cos(\tfrac14P+\tfrac34\varphi)^2,
\]
more accurately still
\[
\varphi=\psi+\frac{ee(\varphi-P)\sqrt[3]{(\cos\frac12(\varphi-P))}.\cos\frac12(\varphi+P)\cos\frac12(\varphi+\psi)}{1-ee\sin\frac12(P+\varphi)^2},
\]
or more accurately still if the denominator [is set] $=1-ee\sin P\sin\varphi$.
\clearpage
\noindent\mbox{}\hfill{\small\textsc{Conformal Projection of the Spheroid onto the Sphere.}}\hfill113\par\bigskip
\begin{center}[8.]\end{center}
If
\[
\begin{aligned}
\operatorname{tang}(45^\circ+\tfrac12\psi)&=\operatorname{tang}(45^\circ+\tfrac12\varphi)\left(\frac{1-e\sin\varphi}{1+e\sin\varphi}\right)^{\frac12e}.\left(\frac{1+e\sin P}{1-e\sin P}\right)^{\frac12e}\\
&=\operatorname{tang}(45^\circ+\tfrac12\varphi).e^u
\end{aligned}
\]
is set, then
\begingroup\footnotesize
\[
\frac{\mathrm d\psi}{\mathrm du}=\left[\cos\psi=\frac{2\operatorname{tang}(45^\circ+\frac12\psi)}{1+\operatorname{tang}(45^\circ+\frac12\psi)^2}=\right]
\frac{\cos\varphi}{1+\sin\varphi.u+\frac12uu+\frac16\sin\varphi.u^3+\frac1{24}u^4+\frac1{120}\sin\varphi.u^5+\cdots}
\]
\par\endgroup
and hence
\[
\psi=\varphi+\cos\varphi.u-\tfrac12\cos\varphi\sin\varphi.uu-\tfrac16\cos\varphi\cos2\varphi.u^3\ldots.
\]
[Furthermore, it follows from
\[
\mathrm du=-\frac{ee\cos\varphi}{1-ee\sin\varphi^2}\mathrm d\varphi=-ee(1+ee\sin\varphi^2+e^4\sin\varphi^4+\cdots)\mathrm d\sin\varphi:\,]
\]
\[
-u=ee(\sin\varphi-\sin P)+\tfrac13e^4(\sin\varphi^3-\sin P^3)+\tfrac15e^6(\sin\varphi^5-\sin P^5)+\cdots
\]
or, putting $\sin\varphi=\sin P+x=\sigma$,
\begingroup\small
\[
\begin{aligned}
u&=-eex-\tfrac13e^4(3x\sigma\sigma-3xx\sigma+x^3)-\tfrac15e^6(5x\sigma^4-10xx\sigma^3+10x^3\sigma\sigma-5x^4\sigma+x^5)\ldots\\
&=-\frac{eex}{1-ee\sigma\sigma}-\ldots.
\end{aligned}
\]
\par\endgroup
[One has]
\begingroup\small
\[
\begin{aligned}
\frac{\cos\varphi}{\cos\psi}\bigg[&=\frac{1+\operatorname{tang}(45^\circ+\frac12\psi)^2}{.\,2\operatorname{tang}(45^\circ+\frac12\psi)}\cos\varphi=\cos iu+\sin\varphi\frac{\sin iu}{i}\bigg]\\
&=1+\sin\varphi.u+\tfrac12uu+\tfrac16\sin\varphi.u^3+\tfrac1{24}u^4+\tfrac1{120}\sin\varphi.u^5+\cdots,
\end{aligned}
\]
\[
\log\frac{\cos\varphi}{\cos\psi}=\sin\varphi.u+\tfrac12\cos\varphi^2.uu-\tfrac13\cos\varphi^2\sin\varphi.u^3-\cdots
\]
\par\endgroup
[and]
\begingroup\footnotesize
\[
\sin\psi=\GaussSquareBracketPair{\frac{\operatorname{tang}(45^\circ+\frac12\psi)^2-1}{\operatorname{tang}(45^\circ+\frac12\psi)^2+1}=\frac{\sin\varphi\cos iu+\dfrac{\sin iu}{i}}{\cos iu+\sin\varphi\dfrac{\sin iu}{i}}=}
\frac{\sin\varphi+u+\frac12\sin\varphi.uu+\frac16u^3+\cdots}{1+\sin\varphi.u+\frac12uu+\frac16\sin\varphi.u^3+\cdots},
\]
\[
\sin\varphi-\sin\psi=\GaussSquareBracketPair{-\frac{\cos\varphi^2\dfrac{\sin iu}{i}}{\cos iu+\sin\varphi\dfrac{\sin iu}{i}}=}-\cos\varphi^2.
\frac{u+\frac16u^3+\frac1{120}u^5+\cdots}{1+\sin\varphi.u+\frac12uu+\frac16\sin\varphi.u^3+\cdots},
\]
\[
\frac{\sin\varphi-\sin\psi}{\cos\psi}=\left[-\cos\varphi\frac{\sin iu}{i}=\right]-\cos\varphi.(u+\tfrac16u^3+\tfrac1{120}u^5+\cdots)\left[=\frac{\mathrm d\log n}{\mathrm d\psi}\right].
\]
\par\endgroup
\vfill\noindent{\footnotesize IX.}\hfill15\hspace*{2cm}\par
\clearpage
\noindent114\hfill\textsc{Posthumous Papers.}\hfill\mbox{}\par\bigskip
\begin{center}[9.]\end{center}
[For $P$ (Göttingen) $=51^\circ31'48''{,}70$ and a flattening of $\dfrac1{302{,}68}$, one obtains:]
\begin{center}
\begingroup\renewcommand{\arraystretch}{1.45}\setlength{\tabcolsep}{12pt}
\begin{tabular}[t]{|r|r|r|}
\hline
$\varphi$&$\psi-\varphi$&\\\hline\hline
$48^\circ\ 0'$&$+36''{,}365$&\raisebox{-.7\baselineskip}{\small$1{,}888$}\\
$10$&$34{,}477$&\raisebox{-.7\baselineskip}{\small$1{,}870$}\\
$20$&$32{,}607$&\raisebox{-.7\baselineskip}{\small$1{,}854$}\\
$30$&$30{,}753$&\raisebox{-.7\baselineskip}{\small$1{,}837$}\\
$40$&$28{,}916$&\raisebox{-.7\baselineskip}{\small$1{,}820$}\\
$50$&$27{,}096$&\raisebox{-.7\baselineskip}{\small$1{,}803$}\\
$49\quad0$&$25{,}293$&\raisebox{-.7\baselineskip}{\small$1{,}786$}\\
$10$&$23{,}507$&\raisebox{-.7\baselineskip}{\small$1{,}770$}\\
$20$&$21{,}737$&\raisebox{-.7\baselineskip}{\small$1{,}752$}\\
$30$&$19{,}985$&\raisebox{-.7\baselineskip}{\small$1{,}736$}\\
$40$&$18{,}249$&\raisebox{-.7\baselineskip}{\small$1{,}718$}\\
$50$&$16{,}531$&\raisebox{-.7\baselineskip}{\small$1{,}702$}\\
$50\quad0$&$14{,}829$&\raisebox{-.7\baselineskip}{\small$1{,}684$}\\
$10$&$13{,}145$&\raisebox{-.7\baselineskip}{\small$1{,}668$}\\
$20$&$11{,}477$&\raisebox{-.7\baselineskip}{\small$1{,}650$}\\
$30$&$9{,}827$&\raisebox{-.7\baselineskip}{\small$1{,}634$}\\
$40$&$8{,}193$&\raisebox{-.7\baselineskip}{\small$1{,}616$}\\
$50$&$6{,}577$&\raisebox{-.7\baselineskip}{\small$1{,}600$}\\
$51\quad0$&$4{,}977$&\raisebox{-.7\baselineskip}{\small$1{,}583$}\\
$10$&$3{,}394$&\raisebox{-.7\baselineskip}{\small$1{,}566$}\\
$20$&$1{,}828$&\raisebox{-.7\baselineskip}{\small$1{,}549$}\\
$30$&$+0{,}279$&\raisebox{-.7\baselineskip}{\small$1{,}532$}\\
$40$&$-1{,}253$&\raisebox{-.7\baselineskip}{\small$1{,}516$}\\
$50$&$2{,}769$&\raisebox{-.7\baselineskip}{\small$1{,}499$}\\
$52\quad0$&$4{,}268$&
\end{tabular}\qquad
\begin{tabular}[t]{|r|r|r|}
\hline
$\varphi$&$\psi-\varphi$&\\\hline\hline
$52^\circ\ 0'$&$-4''{,}268$&\raisebox{-.7\baselineskip}{\small$1{,}482$}\\
$10$&$5{,}750$&\raisebox{-.7\baselineskip}{\small$1{,}465$}\\
$20$&$7{,}215$&\raisebox{-.7\baselineskip}{\small$1{,}449$}\\
$30$&$8{,}664$&\raisebox{-.7\baselineskip}{\small$1{,}432$}\\
$40$&$10{,}096$&\raisebox{-.7\baselineskip}{\small$1{,}415$}\\
$50$&$11{,}511$&\raisebox{-.7\baselineskip}{\small$1{,}399$}\\
$53\quad0$&$12{,}910$&\raisebox{-.7\baselineskip}{\small$1{,}382$}\\
$10$&$14{,}292$&\raisebox{-.7\baselineskip}{\small$1{,}364$}\\
$20$&$15{,}656$&\raisebox{-.7\baselineskip}{\small$1{,}348$}\\
$30$&$17{,}004$&\raisebox{-.7\baselineskip}{\small$1{,}331$}\\
$40$&$18{,}335$&\raisebox{-.7\baselineskip}{\small$1{,}315$}\\
$50$&$19{,}650$&\raisebox{-.7\baselineskip}{\small$1{,}298$}\\
$54\quad0$&$20{,}948$&\raisebox{-.7\baselineskip}{\small$1{,}282$}\\
$10$&$22{,}230$&\raisebox{-.7\baselineskip}{\small$1{,}265$}\\
$20$&$23{,}495$&\raisebox{-.7\baselineskip}{\small$1{,}248$}\\
$30$&$24{,}743$&\raisebox{-.7\baselineskip}{\small$1{,}232$}\\
$40$&$25{,}975$&\raisebox{-.7\baselineskip}{\small$1{,}215$}\\
$50$&$27{,}190$&\raisebox{-.7\baselineskip}{\small$1{,}198$}\\
$55\quad0$&$-28{,}388$&
\end{tabular}
\par\endgroup
\end{center}
\clearpage
\noindent\mbox{}\hfill{\small\textsc{Conformal Projection of the Spheroid onto the Sphere.}}\hfill115\par\bigskip
\begin{center}[10.]\end{center}
The corrections to the computed azimuths [on the sphere] for ellipticity are, [if]
\[
\begin{array}{rccl}
\text{at the first point:}&(\varphi-\psi)\sin\zeta&=&a\\
\text{in the middle:}&\text{»}&=&b\\
\text{at the second point:}&\text{»}&=&c
\end{array}
\]
[and]
\[
s=\text{the entire arc [and $\zeta$ is its southwest azimuth],}
\]
\[
\begin{array}{rl}
\text{at the first point:}&(\tfrac13b+\tfrac16a)s\\
\text{»\quad second\quad »\quad :}&-(\tfrac13b+\tfrac16c)s.
\end{array}
\]
[Thus] the following is added to the \textls[120]{observed} [azimuth]
\[
\begin{array}{rl}
\text{[at the first point:]}&-(\tfrac13b+\tfrac16a)s\\
\text{[ »\quad second\quad »\quad :]}&+(\tfrac13b+\tfrac16c)s.
\end{array}
\]
\begin{center}\rule{.18\textwidth}{.4pt}\par\vspace{12mm}\textsc{Remarks.}\end{center}
\begingroup\small
The preceding formulas for the conformal mapping of the Earth's ellipsoid onto the sphere, and substantially also the formulas in [II] and [III] for mapping the sphere onto the plane, were devised at the same time as the formulas for the conformal representation of the ellipsoid in the plane that were subsequently used in the Hanoverian arc measurement and land survey. They probably date, therefore, from the end of the second decade of the previous century. The mapping of the ellipsoid of revolution onto the sphere as given here, based on Art.~13 of the “Allgemeinen Auflösung der Aufgabe, die Theile einer gegebenen Fläche auf einer andern gegebenen Fläche so abzubilden etc.,” was later superseded by the second mapping given in the “Untersuchungen über Gegenstände der höhern Geodäsie, 1844,” which --- as \textsc{Gauss} himself says (Vol.~IV, p.~262) --- is much better suited to geodetic applications.

The entries [1], [6], [10], as well as [1], p.~117, for the stereographic projection and [1], p.~123, for the Mercator projection, follow one another directly in the same notebook.

The notes [2], [3], [7], [8], and [9] are scattered among formulas for the conformal representation of the Earth's ellipsoid in the plane in another notebook; [4] and [5] are on separate sheets among numerical calculations.
\par\endgroup
\vfill\hfill$15^*$\hspace*{2cm}\par
\clearpage
\noindent116\hfill{\footnotesize\textsc{Remarks.\quad Conformal Projection of the Spheroid onto the Sphere.}}\par\medskip
\begingroup\small
In [3], in the formula for $\dfrac{\mathrm d\log n}{\mathrm d\psi}$, the term multiplied by $(\psi-P)^3$ has $\tfrac23\cos P^2+\tfrac12\sin P^2$ in the original in place of $\tfrac23\cos P^2-\tfrac12\sin P^2$. Accordingly, the last term in the formula for $\log n$ in [4] has $\tfrac16\cos P^2+\tfrac18\sin P^2$ in place of $\tfrac16\cos P^2-\tfrac18\sin P^2$, and the expression for $\gamma$ in [5] has $\tfrac1{30}\cos P^2+\tfrac1{40}\sin P^2$ in place of $\tfrac1{30}\cos P^2-\tfrac1{40}\sin P^2$.

In the formulas for $\psi$ and $\varphi$ in [6], \textsc{Gauss} uses in the last terms, in place of the factor $\tfrac34$, the factor $\tfrac14$. Accordingly (and also because of a small computational error), in the formula $\psi=\varphi-(1)\delta+(2)\delta\delta$ in the example on p.~119, $\log(2)=1{,}86733-10$, as given in the original, has been replaced by $\log(2)=2{,}35885-10$.

The differences $\psi_1$ and $\psi_2$ between the azimuths of the image of the geodesic on the spheroid and the azimuths of the great circle on the sphere, Art.~[10], follow, according to Vol.~IV, p.~278 ff., from the equations:
\[
\begin{aligned}
\psi_1&=\tfrac12l_1s+\tfrac16\lambda ss\ldots\\
\psi_2&=-\tfrac12l_1s-\tfrac13\lambda ss\ldots\\
&=-\tfrac12l_2s+\tfrac16\lambda ss\ldots;
\end{aligned}
\]
$l_1$ and $l_2$ are the values of $l=\dfrac{\mathrm dn}{n\,\mathrm d\psi}\sin\zeta$ at the initial and final points of the geodesic. If $l_m$ is the value of $l$ at the midpoint, then
\[
l_m=l_1+\lambda\frac s2\ldots=l_2-\lambda\frac s2\ldots.
\]
Thus
\[
\psi_1=(\tfrac16l_1+\tfrac13l_m)s\ldots,\qquad\psi_2=-(\tfrac16l_2+\tfrac13l_m)s\ldots.
\]
By [1], p.~108, $l=\dfrac{\sin\varphi-\sin\psi}{\cos\psi}\sin\zeta=(\varphi-\psi)\sin\zeta-\ldots.$\par
\noindent\hfill\textsc{Krüger, Börsch.}\par
\par\endgroup
\vfill\begin{center}\rule{.24\textwidth}{.4pt}\end{center}
\clearpage
% Unnumbered opening at the position of printed 117.
\begin{center}
[II.]\par\bigskip
{\large [STEREOGRAPHIC PROJECTION OF THE SPHERE\\ONTO THE PLANE.]}\par\bigskip
\rule{.1\textwidth}{.4pt}\par\bigskip
[1.]
\end{center}
If the sphere is then mapped onto the plane by \textls[110]{stereographic} projection, the position of the location whose abscissa and ordinate [are] $x$, $y$ is found as follows. Put
\[
\begin{aligned}
2r\operatorname{tang}(45^\circ-\tfrac12P)+x&=A\cos\alpha&\qquad y&=A\sin\alpha\\
2r\operatorname{tang}(45^\circ+\tfrac12P)-x&=B\cos\beta&y&=B\sin\beta.
\end{aligned}
\]
[$x$ is positive toward the south. $P$ is the geographical latitude of the central point. If $n$ and $s$ are the projections of the two poles and $q$ is the projection of any point whose latitude is $\psi$ and whose longitude is $\lambda$, then $nq=A$ and $sq=B$; $\alpha$ and $\beta$ are the angles that $A$ and $B$ make with the $x$-axis.] The longitude of the location is then $=\alpha+\beta$ [with longitude increasing with $y$]; and the latitude is found from
\[
\operatorname{tang}(45^\circ+\tfrac12\psi)=\frac BA\operatorname{tang}(45^\circ-\tfrac12P).
\]
[Furthermore,]
\[
x+y\operatorname{tang}\tfrac12(\alpha-\beta)=2r\operatorname{tang}\tfrac12(P-\psi).
\]
[The] inclination of the location's meridian to the abscissa axis [is] $\alpha-\beta$.
\begin{center}Example.\end{center}
Let $P$ [Marseille] $=43^\circ17'50''{,}1$; [the coordinates of] the island of Planier [are]:
\[
x=+5628{,}56\ \text{Tois.}\qquad y=+5679{,}67\ \text{Tois.}
\]
\clearpage
\noindent118\hfill\textsc{Posthumous Papers.}\hfill\mbox{}\par\medskip
[$r$ was computed using a flattening of $\tfrac1{310}$, cf. p.~68.]
\begin{center}
\begingroup\fontsize{8}{12}\selectfont\setlength{\tabcolsep}{3pt}\renewcommand{\arraystretch}{1.4}
\begin{tabular}{@{}r@{ }r@{\quad}r@{ }r@{\quad}l@{}}
$a\ \text{(in toises)} .\,.\,.\,.\,.\,.$&$6{,}514\,7609$&$y.\,.\,.\,.\,.\,.\,.\,.\,.\,.$&$3{,}754\,3231$&\\
$\sqrt{(1-ee\sin P^2)}.\,.\,.\,.$&$\underline{9{,}999\,3412}$&$A\cos\alpha.\,.\,.\,.\,.\,.$&$6{,}452\,5266$&$\alpha=6'53''{,}258$\\
$r.\,.\,.\,.\,.\,.\,.\,.\,.\,.\,.\,.\,.\,.$&$6{,}515\,4197$&$B\cos\beta.\,.\,.\,.\,.\,.$&$7{,}181\,0749$&$\underline{\beta=1\ 17{,}210}$\\\cline{3-4}
$2r.\,.\,.\,.\,.\,.\,.\,.\,.\,.\,.\,.\,.\,.$&$6{,}816\,4497$&$\dfrac BA.\,.\,.\,.\,.\,.\,.\,.\,.\,.\,.\,.\,.\,.$&$0{,}728\,5474$&$\lambda=8'10''{,}468$\\
$\operatorname{tang}(45^\circ+\tfrac12P).\,.\,.\,.$&$0{,}364\,7863$&$\operatorname{tang}(45^\circ+\tfrac12P).\,.\,.$&$\underline{0{,}364\,7863}$&\\
$2r\operatorname{tang}(45^\circ-\tfrac12P)=$&$2\,829\,199$&$\operatorname{tang}(45^\circ+\tfrac12\psi).\,.\,.$&$0{,}363\,7611$&\\
$2r\operatorname{tang}(45^\circ+\tfrac12P)=$&$15\,178\,749$&$45^\circ+\tfrac12\psi=$&$66^\circ35'57''{,}71$&\\
$x=$&$5628{,}56$&$\psi=$&$43\ 11\ 55{,}42.$&
\end{tabular}\par\endgroup
\end{center}
For the inverse problem one uses the formula
\[
x+iy=2r\frac{\cos\frac12\lambda\sin\frac12(P-\psi)-i\sin\frac12\lambda\cos\frac12(P+\psi)}{\cos\frac12\lambda\cos\frac12(P-\psi)+i\sin\frac12\lambda\sin\frac12(P+\psi)}.
\]
[Here $\lambda$ is positive toward the east and, as before, $x$ is positive toward the south and $y$ toward the west. If the northeast azimuth of the great circle passing through the origin and the location $(\psi,\lambda)$ is $=T$ at the origin and $=T'$ at the location, then
\[
\begin{aligned}
\cos\tfrac12(T'+T)\sin\tfrac12\Delta&=\sin\tfrac12(\psi-P)\cos\tfrac12\lambda\\
\sin\tfrac12(T'+T)\sin\tfrac12\Delta&=\cos\tfrac12(\psi+P)\sin\tfrac12\lambda\\
\cos\tfrac12(T'-T)\cos\tfrac12\Delta&=\cos\tfrac12(\psi-P)\cos\tfrac12\lambda\\
\sin\tfrac12(T'-T)\cos\tfrac12\Delta&=\sin\tfrac12(\psi+P)\sin\tfrac12\lambda.
\end{aligned}
\]
This gives
\[
x+iy=-2re^{iT}\operatorname{tang}\tfrac12\Delta,
\]
hence
\[
x=-2r\operatorname{tang}\frac\Delta2\cos T\qquad y=-2r\operatorname{tang}\frac\Delta2\sin T.\,]
\]
\begin{center}[Example.]\end{center}
\[
\begin{array}{rcl@{\qquad}rl}
\text{Göttingen:}&P=&51^\circ31'55''&\text{Longitude}=&27^\circ35'45''\\
\text{Brocken:}&\varphi=&51\ 48\ 12&&\underline{28\ 17\ 1}\\
&&&\lambda=&+41'16''
\end{array}
\]
[First, by [7], p.~112, with a flattening of $\tfrac1{310}$, one has]
\clearpage
\noindent\mbox{}\hfill{\small\textsc{Stereographic Projection of the Sphere onto the Plane.}}\hfill119\par\medskip
\begin{center}
\begingroup\fontsize{8}{11}\selectfont\setlength{\tabcolsep}{3pt}\renewcommand{\arraystretch}{1.4}
\begin{tabular}{@{}r@{ }r@{\quad}r@{ }r@{\quad}l@{}}
$\sin P.\,.\,.$&$9{,}893\,7368$&$\sin\varphi.\,.\,.$&$9{,}895\,3633$&[One also has:]\\
$e.\,.\,.\,.\,.\,.$&$8{,}904\,4836$&$e.\,.\,.\,.\,.\,.$&$8{,}904\,4836$&$\varphi=P+\delta$\\
$U=$&$3^\circ36'9''{,}755$&$u=$&$3^\circ36'58''{,}484$&$\psi=\varphi-(1)\delta+(2)\delta\delta$\\
$\operatorname{tang}(45^\circ+\tfrac12U).\,.\,.$&$0{,}027\,3261$&&&$(1).\,.\,. 7{,}39837-10$\\
$\operatorname{tang}(45^\circ+\tfrac12u).\,.\,.\,.$&$\underline{0{,}027\,4289}$&&&$(2).\,.\,. 2{,}35885-10$\\
&$0{,}000\,1028$&$1028.\,.\,.\,.\,.\,.\,.\,.$&$3{,}01199$&\\
&&$e.\,.\,.\,.\,.\,.\,.\,.\,.\,.\,.\,.$&$8{,}90448$&$[\delta=977''$\\
&&$\dfrac{206\,265}{\mathrm{Mod.}\,10^7}\cos\varphi.\,.\,.$&$\underline{8{,}46788}$&$\psi=\varphi-2''{,}42]$\\
&&$\varphi-\psi.\,.\,.\,.\,.\,.\,.\,.$&$0{,}38435$&\\
&&\multicolumn{2}{c}{$\varphi-\psi=2''{,}42$}&$\psi=51^\circ48'9''{,}58$
\end{tabular}\par\endgroup
\end{center}
\begin{center}
\begingroup\fontsize{8}{11}\selectfont\setlength{\tabcolsep}{3pt}\renewcommand{\arraystretch}{1.4}
\begin{tabular}{@{}r@{ }r@{\quad}r@{ }r@{\quad}l@{}}
$\sin\tfrac12(\psi-P).\,.\,.\,.$&$7{,}37336$&$\cos\tfrac12(\psi+P).\,.\,.\,.$&$9{,}79255$&$\tfrac12(T'+T)=57^\circ36'\ 0''$\\
$\cos\tfrac12\lambda.\,.\,.\,.\,.\,.\,.\,.$&$9{,}99999$&$\sin\tfrac12\lambda.\,.\,.\,.\,.\,.\,.\,.\,.$&$7{,}77829$&$\tfrac12(T'-T)=\phantom{5}0\ 16\ 11$\\
$\cos\tfrac12(\psi-P).\,.\,.$&$0{,}00000$&$\sin\tfrac12(\psi+P).\,.\,.\,.$&$9{,}89455$&$T=57\ 19\ 49$\\\cline{1-2}\cline{3-4}
$\sin\tfrac12\lambda\cos\tfrac12(\psi+P).\,.$&$7{,}57084$&$\sin\tfrac12\lambda\sin\tfrac12(\psi+P).\,.$&$7{,}67284$&\\
$\cos\tfrac12\lambda\sin\tfrac12(\psi-P).\,.$&$\underline{7{,}37335}$&$\cos\mathord{|}\tfrac12\lambda\cos\tfrac12(\psi-P).\,.$&$\underline{9{,}99999}$&\\
$[\operatorname{tang}\tfrac12(T'+T).\,.\,.\,.$&$0{,}19749$&$\operatorname{tang}\tfrac12(T'-T).\,.\,.\,.\,.$&$7{,}67285]$&\\
$\sin\tfrac12(T'+T).\,.\,.\,.\,.\,.$&$9{,}92651$&$\cos\tfrac12(T'-T).\,.\,.\,.\,.$&$0{,}00000$&\\
$\sin\tfrac12\Delta.\,.\,.\,.\,.\,.\,.\,.\,.\,.\,.$&$7{,}64433$&$\cos\tfrac12\Delta.\,.\,.\,.\,.\,.\,.\,.\,.\,.\,.$&$9{,}99999$&
\end{tabular}\par\endgroup
\end{center}
\begin{center}
\begingroup\small\renewcommand{\arraystretch}{1.15}
\begin{tabular}{r@{ }r@{\qquad}l}
$2a.\,.\,.\,.\,.\,.\,.\,.\,.\,.$&$7{,}105\,6109$&\\
$\cos U.\,.\,.\,.\,.\,.\,.$&$\underline{9{,}999\,1408}$&\\
$2r.\,.\,.\,.\,.\,.\,.\,.\,.\,.$&$7{,}106\,4701$&\\
$\operatorname{tang}\tfrac12\Delta.\,.\,.\,.\,.$&$\underline{7{,}64434}$&\\
$2r\operatorname{tang}\tfrac12\Delta.\,.\,.$&$4{,}75081$&$56339\ \mathrm m=\sqrt{(xx+yy)}$\\
$\cos T.\,.\,.\,.\,.\,.\,.$&$9{,}73223$&$-30411\ \text{»}=x$\\
$\sin T.\,.\,.\,.\,.\,.\,.\,.$&$9{,}92521$&$-47427\ \text{»}=y.$
\end{tabular}\par\endgroup
\end{center}
Twice the auxiliary angle of the denominator [namely $\tfrac12(T'-T)$] gives the convergence of the meridians, or the angle $\alpha-\beta$ above, here $32'22''$.

Finally, the correction to the azimuth of the location $x'$, $y'$, as seen from the location $x$, $y$ and computed along a straight line, is found from the formula:
\[
\operatorname{tang}\mathrm{Corr.}=\frac{x'y-xy'}{4rr+xx'+yy'}\quad\text{exactly}\quad=\frac{\dfrac{2rd}{\operatorname{tang}\frac12\delta}\sin\mathrm{Corr.}}{\dfrac{2rd}{\operatorname{tang}\frac12\delta}\cos\mathrm{Corr.}},
\]
\clearpage
\noindent120\hfill\textsc{Posthumous Papers.}\hfill\mbox{}\par\medskip
where
\[
\begin{array}{rcccl}
d&\text{[is] the}&\text{distance}&\text{in the plane}&\\
\delta&\text{»}&\text{»}&\text{on the sphere}&\text{.}
\end{array}
\]
\bigskip
\begin{center}
[2.]\par\medskip
The stereographic projection onto the horizon of a location\\whose latitude is $B$ [and the] longitude [is] $0$.
\end{center}
Let the latitude of the location to be mapped [be] $\beta$, [and its] longitude $\lambda$. [Then we have]
\[
\begin{aligned}
\text{its projection, northward:}\quad&\frac{\sin\beta\cos B-\cos\beta\sin B\cos\lambda}{1+\sin B\sin\beta+\cos B\cos\beta\cos\lambda}=y\\[6pt]
\text{eastward:}\quad&\frac{\cos\beta\sin\lambda}{1+\sin B\sin\beta+\cos B\cos\beta\cos\lambda}=x
\end{aligned}
\]
[if the diameter of the sphere is taken $=1$].

Or
\[
x=\frac{\cos\beta\sin\lambda}{1+\cos\frac12\lambda^2\cos(\beta-B)-\sin\frac12\lambda^2\cos(\beta+B)}
\]
\[
\begin{aligned}
y&=\frac{\cos\frac12\lambda^2\sin(\beta-B)+\sin\frac12\lambda^2\sin(\beta+B)}{1+\cos\frac12\lambda^2\cos(\beta-B)-\sin\frac12\lambda^2\cos(\beta+B)}\\[6pt]
&=\frac{\dfrac1{\cos B}(\sin B+\sin\beta)}{\mathrm{Denomin.}}-\operatorname{tang}B.
\end{aligned}
\]
The common denominator [is also]
\[
=2\cos\tfrac12\lambda^2\cos\tfrac12(\beta-B)^2+2\sin\tfrac12\lambda^2\sin\tfrac12(\beta+B)^2.
\]
1) The meridians become circles of radius $\dfrac1{\sin\lambda\cos B}$ with center:
\[
\begin{aligned}
y&=-\operatorname{tang}B\\
x&=\frac1{\cos B\operatorname{tang}\lambda}.
\end{aligned}
\]
\bigskip
2) The parallels become circles of radius
\[
\tfrac12\operatorname{cotang}\tfrac12(\beta+B)-\tfrac12\operatorname{tang}\tfrac12(\beta-B)
\]
with center:
\[
\begin{aligned}
y&=\tfrac12\operatorname{cotang}\tfrac12(\beta+B)+\tfrac12\operatorname{tang}\tfrac12(\beta-B)\\
x&=0.
\end{aligned}
\]
\clearpage
\noindent\mbox{}\hfill{\small\textsc{Stereographic Projection of the Sphere onto the Plane.}}\hfill121\par\bigskip
\begin{center}[3.]\end{center}
[Let]
\[
\begin{gathered}
\left.\begin{array}{ll}
2p\, .\, .\, .\, .&\text{polar distance}\\
\lambda\, .\, .\, .\, .\, .&\text{longitude}
\end{array}\right\}\ \text{[of any point on the sphere]}\\[5pt]
2q\, .\, .\, .\, .\ \text{polar distance of the central point;}
\end{gathered}
\]
[furthermore, let the $x$-axis in the plane be the image of the initial meridian, with its positive direction toward the south. The ordinate $y$ has the same sign as $\lambda$.

If the diameter of the sphere is $=1$, then]
\[
x+iy=\frac{\operatorname{tang}p(\cos\lambda+i\sin\lambda)-\operatorname{tang}q}{1+\operatorname{tang}q\operatorname{tang}p(\cos\lambda+i\sin\lambda)}.
\]
[Put]
\[
\begin{aligned}
\cos\lambda\operatorname{tang}p-\operatorname{tang}q&=a\cos A\\
\sin\lambda\operatorname{tang}p&=a\sin A\\
1+\operatorname{tang}q\operatorname{tang}p\cos\lambda&=b\cos B\\
\operatorname{tang}q\operatorname{tang}p\sin\lambda&=b\sin B.
\end{aligned}
\]
[This gives]
\[
x+iy=\frac ab\{\cos(A-B)+i\sin(A-B)\}.
\]
More simply still
\[
x+iy=\frac{\sin(p-q)\cos\frac12\lambda+i\sin(p+q)\sin\frac12\lambda}{\cos(p-q)\cos\frac12\lambda-i\cos(p+q)\sin\frac12\lambda}.
\]
\begin{center}\rule{.18\textwidth}{.4pt}\par\vspace{10mm}\textsc{Remarks.}\end{center}
\begingroup\small
As already mentioned on p.~115, note [1] is in a notebook; in the numerical examples, the labels have been added and a few small computational inaccuracies removed. The azimuth reduction at the end of this note can be obtained as follows. Let $x$ be positive toward the south and let $y$ have the same sign as $\lambda$. Then
\[
x=-2r\operatorname{tang}\frac\Delta2\cos T=-\rho\cos T,\qquad y=2r\operatorname{tang}\frac\Delta2\sin T=\rho\sin T,
\]
where $\Delta$ is thus the great-circle arc that joins the point to be mapped to the central point, and
\par\endgroup
\vfill\noindent{\footnotesize IX.}\hfill16\hspace*{2cm}\par
\endgroup
